\documentclass{stex}

\libusepackage{naproche}
\libusepackage{copyright}
\libinput{preamble}
\libinput[libraries/meta]{preamble}

\title{The Ontology of \Naproche}
\author{Marcel Schütz}
\date{2025}

\begin{document}

\usemodule[libraries/meta]{preliminaries.ftl}

\maketitle

\begin{abstract}
  This document provides an overview of the object-level ontology of \Naproche.
  Moreover, it introduces common (symbolic and verbal) macros for the notions
  provided by this ontology which extend the core language of \ForTheL.
\end{abstract}

\ifstexhtml\else\tableofcontents\fi


\section{Introduction}\label{sec:introduction}

This document describes the notions and axioms hard-coded in \Naproche.
Furthermore, it presents some additional symbolic and verbal macros that extend
the core language of \ForTheL.
These macros are provided by the library \path{libraries/meta} which contains a
module \path{preliminaries.ftl} that exports all these macros.
Every library module (excluding those in the library \path{meta}) is expected to
import the module \path{preliminaries.ftl}.
So by importing any library module to your formalization, you automatically have
access to all macros described in this document.


\section{Propositions}\label{sec:propositions}

\ForTheL is essentially a surface language for (unsorted) first-order logic.
We therefore provide some symbolic notations for logical connectives and
quantifiers that are supported by \ForTheL:


\subsection{Truth}\label{sec:truth}

\begin{sparagraph}[style=symdoc,for=TRUE]
  \begin{signature*}
    $\fakeemph{\TRUE}$ is a formula.
  \end{signature*}

  \begin{axiom*}
    $\TRUE$ is true.
  \end{axiom*}
\end{sparagraph}


\subsection{Falsity}\label{sec:falsity}

\begin{sparagraph}[style=symdoc,for=FALSE]
  \begin{signature*}
    $\fakeemph{\FALSE}$ is a formula.
  \end{signature*}

  \begin{axiom*}
    $\FALSE$ is false.
  \end{axiom*}
\end{sparagraph}


\subsection{Negation}\label{sec:negation}

\begin{sparagraph}[style=symdoc,for=NOT]
  \begin{signature*}
    Let $\varphi$ be a formula.
    $\fakeemph{\NOT \varphi}$ is a formula.
  \end{signature*}

  \begin{axiom*}
    Let $\varphi$ be a formula.
    $\NOT\varphi$ is true iff $\varphi$ is false.
  \end{axiom*}
\end{sparagraph}


\subsection{Conjunction}\label{sec:conjunction}

\begin{sparagraph}[style=symdoc,for=AND]
  \begin{signature*}
    Let $\varphi, \psi$ be formulas.
    $\fakeemph{\varphi \AND \psi}$ is a formula.
  \end{signature*}

  \begin{axiom*}
    Let $\varphi, \psi$ be formulas.
    $\varphi \AND \psi$ is true iff $\varphi$ is true and $\psi$ is true.
  \end{axiom*}
\end{sparagraph}


\subsection{Disjunction}\label{sec:disjunction}

\begin{sparagraph}[style=symdoc,for=OR]
  \begin{signature*}
    Let $\varphi, \psi$ be formulas.
    $\fakeemph{\varphi \OR \psi}$ is a formula.
  \end{signature*}

  \begin{axiom*}
    Let $\varphi, \psi$ be formulas.
    $\varphi \OR \psi$ is true iff $\varphi$ is true or $\psi$ is true.
  \end{axiom*}
\end{sparagraph}


\subsection{Implication}\label{sec:implication}

\begin{sparagraph}[style=symdoc,for=IMPLIES]
  \begin{signature*}
    Let $\varphi, \psi$ be formulas.
    $\fakeemph{\varphi \IMPLIES \psi}$ is a formula.
  \end{signature*}

  \begin{axiom*}
    Let $\varphi, \psi$ be formulas.
    $\varphi \IMPLIES \psi$ is true iff $\varphi$ is true implies that $\psi$ is true.
  \end{axiom*}
\end{sparagraph}


\subsection{Equivalence}\label{sec:equivalence}

\begin{sparagraph}[style=symdoc,for=IFF]
  \begin{signature*}
    Let $\varphi, \psi$ be formulas.
    $\fakeemph{\varphi \IFF \psi}$ is a formula.
  \end{signature*}

  \begin{axiom*}
    Let $\varphi, \psi$ be formulas.
    $\varphi \IFF \psi$ is true iff $\varphi$ is true if and only if $\psi$ is true.
  \end{axiom*}
\end{sparagraph}


\subsection{Universal Quantification}\label{sec:universal-quantification}

\begin{sparagraph}[style=symdoc,for=FORALL]
  \begin{signature*}
    Let $x$ be a variable and $\tau$ be a term and $R$ be a binary relation and $\varphi$ be a formula.
    $\fakeemph{\FORALL} x \mathrel{R} \tau : \varphi$ is a formula.
  \end{signature*}

  \begin{axiom*}
    Let $x$ be a variable and $\tau$ be a term and $R$ be a binary relation and $\varphi$ be a formula.
    $\FORALL x \mathrel{R} \tau : \varphi$ is true iff $\varphi[\tau/x]$ is true for all terms $\sigma$ such that $\sigma \mathrel{R} \tau$.
  \end{axiom*}
\end{sparagraph}


\subsection{Existential Quantification}\label{sec:existential-quantification}

\begin{sparagraph}[style=symdoc,for=EXISTS]
  \begin{signature*}
    Let $x$ be a variable and $\tau$ be a term and $R$ be a binary relation and $\varphi$ be a formula.
    $\fakeemph{\EXISTS} x \mathrel{R} \tau : \varphi$ is a formula.
  \end{signature*}

  \begin{axiom*}
    Let $x$ be a variable and $\tau$ be a term and $R$ be a binary relation and $\varphi$ be a formula.
    $\EXISTS x \mathrel{R} \tau : \varphi$ is true iff $\varphi[\tau/x]$ is true for some term $\sigma$ such that $\sigma \mathrel{R} \tau$.
  \end{axiom*}
\end{sparagraph}


\section{Notions}\label{sec:notions}

\subsection{Objects}\label{sec:objects}

\emph{Objects} are mathematical entities that are ``small'' enough to be
contained in classes.

\begin{fakeforthel}
  \begin{signature*}
    A \emph{mathematical object} is a notion.
    Let an \emph{object} stand for a mathematical object.
    Let an \emph{element} stand for a mathematical object.
  \end{signature*}
\end{fakeforthel}

\subsection{Ordered Pairs}\label{sec:ordered-pairs}

\begin{fakeforthel}
  \begin{signature*}
    Let $a,b$ be objects.
    $\emph{(a,b)}$ is an object.
    Let the \emph{ordered pair of $a$ and $b$} stand for $(a,b)$.
  \end{signature*}
\end{fakeforthel}

\begin{fakeforthel}
  \begin{axiom*}[forthel,title=Ordered Pair Extensionality]
    Let $a,a',b,b'$ be objects.
    If $(a, b) \Eq (a', b')$ then $a \Eq a'$ and $b \Eq b'$.
  \end{axiom*}
\end{fakeforthel}


\subsection{Classes and Sets}\label{sec:classes-and-sets}

\begin{fakeforthel}
  \begin{signature*}
    A \emph{class} is a notion.
    Let a \emph{collection} stand for a class.
  \end{signature*}
\end{fakeforthel}

\begin{fakeforthel}
  \begin{axiom*}[title=Class Extensionality]
    Let $A,B$ be classes.
    If every element of $A$ is an element of $B$ and every element of $B$ is an element of $A$ then $A \Eq B$.
  \end{axiom*}
\end{fakeforthel}

Sets are considered to be classes that are ``small'' enough to be contained in
other classes:

\begin{fakeforthel}
  \begin{definition*}
    A \emph{set} is a class that is an object.
  \end{definition*}
\end{fakeforthel}

\begin{sparagraph}[style=symdoc,for={Elem,NotElem}]
  \begin{fakeforthel}
    \begin{signature*}
      Let $A$ be a class.
      A \emph{element of $A$} is an object.
      Let a \emph{member of $X$} stand for an element of $X$.
      Let $x$ \emph{belongs to $X$} stand for $x$ is an element of $X$.
      Let $X$ \emph{contains $x$} stand for $x$ is an element of $X$.
      Let $x$ is \emph{contained in $X$} stand for $x$ is an element of $X$.
      Let $x$ \emph{lies in $X$} stand for $x$ is an element of $X$.
      Let $x$ is \emph{in $X$} stand for $x$ is an element of $X$.
      Let $\emph{a \Elem A}$ stand for $a$ is an element of $A$.
      Let $\emph{a \NotElem A}$ stand for $\NOT a \Elem A$.
    \end{signature*}
  \end{fakeforthel}
\end{sparagraph}


\subsection{Comprehension Terms}\label{sec:comprehension-terms}

In \ForTheL we have three types of comprehension terms:
\begin{itemize}
  \item ``Enumeration classes'', i.e. terms of the form
  $\EClass{a_1,\dots,a_n}$, that represent the class that contains precisely the
  objects $a_1,\dots,a_n$.
  Typical examples are $\EClass{4,7,11}$ (i.e. the class consisting of the
  numbers $4$, $7$ and $11$) or $\EClass{x,y+1,z-x}$ (i.e. the class consisting
  of the elements $x$, $y+1$ and $z-x$ for some fixed objects $x$, $y$ and $z$).
  \item ``Comprehension classes'', i.e. terms of the form
  $\CClass{\tau}{\varphi}$, that represent the class of all objects of the form
  $\tau$ such that the components of $\tau$ satisfy the condition $\varphi$.
  Typical examples are $\CClass{p}{p \text{ is an odd prime number}}$ (i.e the
  class of all odd prime numbers) or $\CClass{(p,n)}{n \in \mathbb{N}
  \text{ and } p \text{ is a prime divisor of } n}$ (i.e. the class of all pairs
  $(p,n)$ of natural numbers such that $p$ is a prime divisor of $n$).
  \item ``Separation classes'', i.e. terms of the form
  $\SClass{\tau}{A}{\varphi}$, that represent the subclass of a class $A$ of all
  objects of the form $\tau$ such that the components of $\tau$ satisfy the
  condition $\varphi$.
  Typical examples are $\SClass{p}{\mathbb{P}}{p \text{ is odd}}$ (i.e the
  class of all odd prime numbers) or $\SClass{(p,n)}{\mathbb{N}\times\mathbb{N}}
  {p \text{ is a prime divisor of } n}$ (i.e. the class of all pairs
  $(p,n)$ of natural numbers such that $p$ is a prime divisor of $n$).
\end{itemize}

\begin{sparagraph}[style=symdoc,for=EClass]
  \begin{signature*}
    Let $a_1, \dots, a_n$ be objects.
    $\fakeemph{\EClass{a1, \dots, a_n}}$ is a class.
  \end{signature*}

  \begin{axiom*}
    Let $a_1, \dots, a_n$ be objects.
    Then $a \Elem \EClass{a_1, \dots, a_n}$ iff $a \Eq a_1$ or \dots{} or $a \Eq a_n$.
  \end{axiom*}
\end{sparagraph}

\begin{sparagraph}[style=symdoc,for=CClass]
  \begin{signature*}
    Let $\tau$ be a term and $\varphi$ be a formula.
    $\fakeemph{\CClass{\tau}{\varphi}}$ is a class.
  \end{signature*}

  \begin{axiom*}
    Let $\tau[x_1, \dots, x_n]$ be a term and $\varphi[x_1, \dots, x_n]$ be a formula.
    Then $a \Elem \CClass{\tau}{\varphi}$ iff
    $a$ is an object and there exist $a_1, \dots, a_n$ such that $\varphi[a_1, \dots, a_n]$ and $a \Eq \tau[a_1, \dots, a_n]$.
  \end{axiom*}

  \begin{proposition*}
    Let $\varphi[x]$ be a formula.
    Then $a \Elem \CClass{x}{\varphi}$ iff $a$ is an object and $\varphi[a]$.
  \end{proposition*}
\end{sparagraph}

\begin{sparagraph}[style=symdoc,for=SClass]
  \begin{signature*}
    Let $\tau$ be a term, $A$ be a class and $\varphi$ be a formula.
    $\fakeemph{\SClass{\tau}{A}{\varphi}}$ is a class.
  \end{signature*}

  \begin{axiom*}
    Let $\tau[x_1, \dots, x_n]$ be a term, $A$ be a class and $\varphi[x_1,\dots,x_n]$ be a formula.
    Then $a \Elem \SClass{\tau}{A}{\varphi}$ iff
    $a \Elem A$ and there exist $a_1, \dots, a_n$ such that $\varphi[a_1, \dots, a_n]$ and $a \Eq \tau[a_1, \dots, a_n]$.
  \end{axiom*}

  \begin{proposition*}
    Let $A$ be a class and $\varphi[x]$ be a formula.
    Then $a \Elem \SClass{x}{A}{\varphi}$ iff
    $a \Elem A$ and $\varphi[a]$.
  \end{proposition*}

  \begin{axiom*}
    Let $\tau$ be a term, $A$ be a class and $\varphi$ be a formula.
    If $A$ is a set then $\SClass{\tau}{A}{\varphi}$ is a set.
  \end{axiom*}
\end{sparagraph}


\subsection{Maps and Functions}\label{sec:maps-and-functions}

\begin{fakeforthel}
  \begin{signature*}
    A \emph{map} is a notion.
  \end{signature*}
\end{fakeforthel}

\begin{fakeforthel}
  \begin{axiom*}[title=Map Extensionality]
    Let $f,g$ be maps.
    If $\Dom(f) \Eq \Dom(g)$ and $f(a) \Eq g(a)$ for all $a \Elem \Dom(f)$ then $g \Eq g$.
  \end{axiom*}
\end{fakeforthel}

Similar to the notion of sets which are considered to be ``small'' classes, we
consider \emph{functions} to be ``small'' maps.

\begin{fakeforthel}
  \begin{definition*}
    A \emph{function} is a map that is an object.
  \end{definition*}
\end{fakeforthel}

\begin{forthel}
  \begin{axiom*}
    Let $f$ be a map.
    $f$ is a function iff $\Dom(f)$ is a set.
  \end{axiom*}
\end{forthel}

\begin{sparagraph}[style=symdoc,for=Dom]
  \begin{fakeforthel}
    \begin{signature*}
      Let $f$ be a map.
      $\emph{\Dom(f)}$ is a class.
      Let the \emph{domain of $f$} stand for $\Dom(f)$.
      Let $f$ is \emph{defined on $X$} stand for $\Dom(f) \Eq X$.
    \end{signature*}
  \end{fakeforthel}
\end{sparagraph}

\begin{fakeforthel}
  \begin{signature*}
    Let $f$ be a map and $a \Elem \Dom(f)$.
    $\emph{f(a)}$ is an object.
    Let the \emph{value of $f$ at $x$} stand for $f(x)$.
    Let $\emph{f(x,y)}$ stand for $f((x,y))$.
  \end{signature*}
\end{fakeforthel}


\subsection{$\lambda$-Terms}\label{sec:lambda-terms}

In low-level map definitions, \ForTheL allows to define maps via
$\lambda$-terms:

\begin{sparagraph}[style=symdoc,for=Map]
  \begin{signature*}
    Let $x$ be a variable, $A$ be a class and $\tau$ be a term.
    $\fakeemph{\Map{x}{A}\tau}$ is a map.
  \end{signature*}

  \begin{axiom*}
    Let $x$ be a variable, $A$ be a class and $\tau[x]$ be a term.
    Then $\Dom(\Map{x}{A}\tau) \Eq A$.
  \end{axiom*}

  \begin{axiom*}
    Let $x$ be a variable, $A$ be a class and $\tau[x]$ be a term and $a \Elem A$.
    Then $(\Map{x}{A}\tau)(a) \Eq \tau[a]$.
  \end{axiom*}
\end{sparagraph}


\section{Relations}\label{sec:relations}

\begin{sparagraph}[style=symdoc,for=Eq]
  \begin{signature*}
    Let $x, y$ be terms.
    $\fakeemph{x \Eq y}$ is a relation.
    Let $x$ is \emph{equal to $y$} stand for $x \Eq y$.
    Let $x$ and $y$ \emph{agree} stand for $x \Eq y$.
    Let $x$ \emph{agrees with $y$} stand for $x \Eq y$.
  \end{signature*}
\end{sparagraph}

\begin{sparagraph}[style=symdoc,for=NotEq]
  \begin{definition*}
    Let $x, y$ be terms.
    $\fakeemph{x \NotEq y}$ iff $\NOT x \Eq y$.
    Let $x$ and $y$ are \emph{distinct} stand for $x \NotEq y$.
  \end{definition*}
\end{sparagraph}

\Naproche has a hard-coded induction proof tactic (cf. section
\ref{sec:misc-axioms}) that relies on an ``induction ordering'' $\ILess$.
An important special case of this tactic is the following induction principle:
If, for all $x$, we can show that a predicate $\varphi[y]$ holds for every
$y \ILess x$, then $\varphi[x]$ holds for every $x$.
Typical instances of this principle are the (strong) induction principle from
natural number arithmetics (with $\ILess$ being the $<$-relation) or the
$\in$-induction principle from set theory (with $\ILess$ being the
$\in$-relation).

\begin{sparagraph}[style=symdoc,for=ILess]
  \begin{signature*}
    Let $x, y$ be terms.
    $\fakeemph{x \ILess y}$ is a relation.
    Let $x$ is \emph{inductively less than $y$} stand for $x \ILess y$.
  \end{signature*}

  \begin{axiom*}
    Let $\tau[x_1, \dots, x_n]$ be a term and
    $\psi_1[x_1]$, $\psi_2[x_1,x_2]$, \dots, $\psi_n[x_1, \dots, x_n]$ and $\varphi[x_1, \dots, x_n]$ be formulas.
    Assume that the following holds for all $a_1, \dots, a_n$ with $\psi_i[a_1, \dots, a_i]$ for all $i \Elem \EClass{1, \dots, n}$:
    If $\tau[y_1, \dots, y_n] \ILess \tau[a_1, \dots, a_n]$
    for all $y_1, \dots, y_n$ with $\psi_i[y_1, \dots, y_i]$ for all $i \Elem \EClass{1, \dots, n}$
    then $\varphi[a_1, \dots, a_n]$.
    Then $\varphi[a_1, \dots, a_n]$ for all $a_1, \dots, a_n$ with $\psi_i[a_1, \dots, a_i]$ for all $i \Elem \EClass{1, \dots, n}$.
  \end{axiom*}

  \begin{proposition*}
    Let $A$ be a class, and $\varphi[x]$ be a formula.
    Assume that the following holds for all $a \Elem A$:
    If $\varphi[b]$ for for all $b \Elem A$ with $b \ILess a$ then $\varphi[a]$.
    Then $\varphi[a]$ for all $a \Elem A$.
  \end{proposition*}
\end{sparagraph}


\printlicense[CcByNcSa]{Marcel Schütz (2025)}

\end{document}
